\documentclass{article}

\usepackage{amsmath} % math stuff
\usepackage{amssymb} % math stuff
\usepackage{array} % equations and stuff
\usepackage{bm} % bold math
%\usepackage{booktabs} % extra table rule options
%\usepackage{caption} % suppressed table numbering; incompatible with revtex, and longtable, I think
\usepackage{comment} % comment environment
%\usepackage{enumitem} % customization of enumeration, itemize, and description
\usepackage[T1]{fontenc} % font encoding for special characters, must also use scalable font package
\usepackage[margin=0.8in]{geometry} % paper sizes and margins (but be careful not to mess up pre-defined pages)
\usepackage{graphicx} % for graphics
%\usepackage{helvet} % default font is the helvetica postscript font
\usepackage[utf8]{inputenc} % special characters in tex input
\usepackage{layouts} % print units like widths
\usepackage{lipsum} % lorem ipsum filler text
\usepackage{lmodern} % scalable font?
\usepackage{longtable} % multi-page tables
\usepackage{makecell} % specify line-breaks in table cells
\usepackage{mathrsfs} % math script font
\usepackage{mhchem} % easier chemical formula
\usepackage{microtype} % allows disabling of ligatures
\usepackage{multicol} % multicolumns
%\usepackage{newcent} % new century schoolbook font
\usepackage{nicefrac}
\usepackage{numprint} % print and format (large) numbers
\usepackage{parskip} % removes paragraph indentation, and adjusts paragraph skip, as well as list items
\usepackage{pdfpages} % add pdf files as pages
%\usepackage{setspace} % adjust text spacing and indents
\usepackage{siunitx} % decimal alignment, and degree symbol
\usepackage{subfigure} % divided figures
%\usepackage{tabu} % extra table options
\usepackage{textcomp} % symbols
\usepackage{threeparttablex} % better footnotes with longtable
\usepackage{titling} % title placement
\usepackage{ulem} % strikethrough text
\usepackage{upgreek} % upright Greek letters
%\usepackage{url} % superceded by hyperref
\usepackage{verbatim} % verbatim environment
\usepackage{xcolor} % colors and color boxes
\usepackage{xspace} % commands that don't eat up white space
\usepackage{hyperref} % links and page setup; should always come last

\hypersetup{
 bookmarks=true,
 colorlinks=true,
 citecolor=blue,
 linkcolor=blue,
 urlcolor=blue,
 pdfstartview={XYZ null null 1.0} % default open view is 100%
}

\DisableLigatures[f,t]{encoding = T1} % disable ff, fi, fl, tt ligatures; without options, it also disables -- = endash
\renewcommand{\arraystretch}{1.0} % extra vertical (and horizontal?) space in tables

% define centered, left- and right-aligned columns with specified widths
\newcommand{\PreserveBackslash}[1]{\let\temp=\\#1\let\\=\temp}
\newcolumntype{C}[1]{>{\PreserveBackslash\centering}p{#1}}
\newcolumntype{L}[1]{>{\PreserveBackslash\raggedright}p{#1}}
\newcolumntype{R}[1]{>{\PreserveBackslash\raggedleft}p{#1}}

\begin{document}

\pagestyle{empty} % don't number pages

% custom title
\begin{center}
{\LARGE Express Riddler}

\vspace{0.15in}

{\Large 8 April 2022}
\end{center}


\section*{Riddle:}

In tic-tac-toe, you can rate a square's ``real estate value'' by how many possible three-in-a-rows pass through it.
Thus, the center scores 4, the corners score 3 and the edges score 2.

That means the center is the ``most valuable'' real estate, while the edges are the ``least valuable''.

Now, let's try three-dimensional tic-tac-toe, played on a 3-by-3-by-3 board.
Which positions are the most and least valuable, and how much are they worth?

\textit{Extra credit:} What about four-dimensional tic-tac-toe?


\section*{Solution:}

There are four types of spaces in three-dimensional tic-tac-toe:
one center, six face centers, twelve edge centers, and eight vertices, for a total of $3^{3}=27$ spaces.
The center space has three orthogonal three-in-a-rows, six along a single diagonal, and four along two diagonals, for a total score of 13.
Each face-center space has three orthogonal three-in-a-rows and two along a single diagonal, for a total score of 5.
Each edge-center space has three orthogonal three-in-a-rows and one along a single diagonal, for a total score of 4.
Each vertex space has three orthogonal three-in-a-rows, three along a single diagonal, and one along two diagonals, for a total score of 7.

There are five types of spaces in four-dimensional tic-tac-toe:
one center, eight ``cube'' centers, 24 face centers, 32 edge centers, and 16 vertices, for a total of $3^{4}=81$ spaces.
The center space has four orthogonal three-in-a-rows, 12 along a single diagonal, 16 along two diagonals, and eight along three diagonals, for a total score of 40.
Each cube-center space has four orthogonal three-in-a-rows, six along a single diagonal, and four along two diagonals, for a total score of 14.
Each face-center space has four orthogonal three-in-a-rows and three along a single diagonal for a total score of 7.
Each edge-center space has four orthogonal three-in-a-rows, three along a single diagonal, and one along two diagonals, for a total score of 8.
Each vertex space has four orthogonal three-in-a-rows, six along a single diagonal, four along two diagonals, and one along three diagonals, for a total score of 15.


\end{document}